\section{What the Fathers Actually Say}

The patristic evidence is rich but does not supply the proposed time-travel
formula. Its recurrent centers are reality of Christ's flesh, ecclesial
offering, memorial, heavenly worship, and participation in the one sacrifice.

Justin Martyr describes bread, wine, thanksgiving, communion, and distribution,
and says that Christians receive the food as the flesh and blood of the
incarnate Jesus (\emph{First Apology} 65--67). Irenaeus argues from the
Eucharist against contempt for creation: earthly bread receives God's
invocation and becomes Eucharist, and communicants' bodies are destined for
resurrection (\emph{Against Heresies} IV.18.4--5). Both witnesses join
sacramental realism to creation and resurrection.

Ambrose links Christ's effective word to the institution narrative and gives a
Western sequence of anamnesis, offering, the sacrifices of Abel, Abraham, and
Melchizedek, and the heavenly altar (\emph{De sacramentis} IV.5.21--23;
IV.6.27). Augustine calls the whole redeemed city a universal sacrifice
offered through the great High Priest and says that the Church celebrates this
in the sacrament of the altar (\emph{City of God} X.6).

These four edition-identified witnesses are representative rather than an
exhaustive patristic survey. They confess reality, offering, memorial, creation,
resurrection, and ecclesial participation. Their bold sacramental language
does not require literal relocation to a future or past civil date.

